% =====================================================================
% Capítulo 2. Objetivos
% Estado: 2.1 y 2.2 redactadas; 2.3 (criterios de cumplimiento) pendiente.
% =====================================================================

\chapter{Objetivos}\label{cap:objetivos}

Este capítulo formula el objetivo general del trabajo y lo descompone en
objetivos específicos. Los objetivos se enuncian ajustados al alcance
delimitado en la sección~\ref{sec:problema}: un único horizonte de previsión,
el del día siguiente, un periodo de estudio de 2017 a 2024 y un conjunto de
modelos reducido a un representante de cada familia. El capítulo termina con
los criterios que permitirán considerar cumplido cada objetivo, que se retoman
uno por uno en las conclusiones.

\section{Objetivo general}\label{sec:objetivo-general}

El objetivo general del trabajo es comparar, bajo un protocolo experimental
equitativo, un modelo estadístico clásico de series temporales con
estacionalidad múltiple frente a modelos de aprendizaje automático y profundo
en la previsión del día siguiente de la demanda eléctrica horaria del sistema
peninsular español, evaluando tanto la precisión puntual de las previsiones
como la calibración de los intervalos de predicción que produce cada familia
de modelos. La previsión se emite a las 12:00 del día $D$ para las 24 horas
del día $D+1$, y la comparación se realiza sobre el periodo 2017-2024, con el
año 2024 como periodo de prueba.

\section{Objetivos específicos}\label{sec:objetivos-especificos}

El objetivo general se concreta en los nueve objetivos específicos siguientes,
ordenados según la secuencia de trabajo: construcción de los datos, análisis
de la serie, modelización, protocolo de comparación, evaluación y
reproducibilidad.

\begin{enumerate}
	
	\item[\textbf{OE1.}] Construir un conjunto de datos reproducible para el
	periodo 2017-2024 que integre la demanda horaria peninsular publicada por
	Red Eléctrica de España a través de ESIOS, la temperatura horaria de
	reanálisis ponderada por la población de cada provincia y el calendario
	laboral estatal y autonómico, documentando el tratamiento de los huecos, de
	los cambios de hora y de los valores anómalos, así como la exclusión de los
	meses de marzo a junio de 2020.
	
	\item[\textbf{OE2.}] Caracterizar estadísticamente la serie de demanda
	mediante la descomposición en tendencia y estacionalidades de 24, 168 y 8760
	horas, el análisis espectral, los contrastes de estacionariedad ADF y KPSS y
	la modelización de la relación no lineal con la temperatura a través de
	grados-día de calefacción y de refrigeración.
	
	\item[\textbf{OE3.}] Especificar, estimar y diagnosticar los modelos
	clásicos de referencia: un método ingenuo estacional, que actúa como suelo de
	rendimiento, y un modelo SARIMAX con términos de Fourier y regresores
	exógenos, validando la adecuación de sus residuos con el contraste de
	Ljung-Box.
	
	\item[\textbf{OE4.}] Implementar los modelos de aprendizaje automático y
	profundo: un modelo de \textit{gradient boosting} (LightGBM) sobre variables
	de retardo y de calendario, y una red neuronal de previsión multihorizonte
	(N-HiTS), ambos capaces de incorporar la temperatura como covariable conocida
	a futuro.
	
	\item[\textbf{OE5.}] Definir y aplicar un protocolo de comparación
	equitativa, basado en validación con origen móvil, idénticas variables
	exógenas, idéntico horizonte de previsión e idéntico presupuesto de búsqueda
	de hiperparámetros para todas las familias, y documentarlo de forma explícita
	como aportación metodológica del trabajo.
	
	\item[\textbf{OE6.}] Cuantificar la precisión puntual de cada modelo mediante
	MAE, RMSE, MAPE y MASE, y contrastar la significación estadística de las
	diferencias entre modelos con el test de Diebold-Mariano, sin extraer
	conclusiones de diferencias que no resulten significativas.
	
	\item[\textbf{OE7.}] Generar intervalos de predicción para cada modelo, en
	primer lugar con el procedimiento propio de cada uno (intervalos analíticos
	del SARIMAX, regresión cuantílica en LightGBM y N-HiTS, y cuantiles empíricos
	del error de validación en el método ingenuo) y, en segundo lugar, mediante
	recalibración con inferencia conforme adaptativa (ACI) aplicada de manera
	uniforme a todos ellos; y evaluar su calidad con la pérdida pinball, el CRPS,
	la cobertura empírica frente a la nominal, la anchura media y los diagramas
	de fiabilidad.
	
	\item[\textbf{OE8.}] Desglosar los resultados de precisión y de calibración
	por hora del día y por tipo de día (laborable, festivo y puente) para
	identificar en qué regímenes concretos resulta preferible cada familia de
	modelos, en lugar de comparar únicamente promedios agregados.
	
	\item[\textbf{OE9.}] Publicar el código, las semillas y las particiones
	temporales en un repositorio público que garantice la reproducibilidad
	íntegra de los experimentos.
	
\end{enumerate}

Respecto a la propuesta aprobada, estos objetivos mantienen su estructura y
su numeración lógica, con las siguientes acotaciones: el horizonte se fija
en el día siguiente, el periodo se establece en 2017-2024, los modelos de
suavizado exponencial, TBATS y \textit{Temporal Fusion Transformer} se tratan
en el marco teórico pero no se implementan, la recalibración conforme se
limita al método ACI y la lectura económica del error de previsión queda como
línea de trabajo futuro. La justificación de cada acotación se expone en el
marco metodológico.

\section{Criterios de cumplimiento}\label{sec:criterios}

La tabla~\ref{tab:criterios} establece, para cada objetivo específico, la
evidencia que permitirá considerarlo cumplido. Los criterios se formulan en
términos de resultados comprobables en la memoria o en el repositorio, y no
en términos del rendimiento de los modelos: un objetivo de comparación se
cumple cuando la comparación se ha realizado correctamente, sea cual sea el
modelo que resulte preferible.

\begin{table}[htbp]
	\centering
	\caption{Criterios de cumplimiento de los objetivos específicos}
	\label{tab:criterios}
	\small
	\begin{tabularx}{\textwidth}{l X}
		\hline
		\textbf{Objetivo} & \textbf{Se considera cumplido si} \\
		\hline
		OE1 & El conjunto de datos se genera íntegramente con el código del repositorio a partir de las fuentes públicas, y la memoria documenta el número de huecos tratados, el criterio aplicado en los cambios de hora y los valores anómalos detectados. \\
		OE2 & La memoria presenta la descomposición de la serie, el periodograma con los picos en los tres periodos estacionales, los resultados de los contrastes ADF y KPSS y la curva demanda-temperatura ajustada con grados-día. \\
		OE3 & Ambos modelos clásicos producen previsiones para todo el periodo de prueba, y el SARIMAX incluye la especificación final (órdenes, número de armónicos y exógenas) y el resultado del contraste de Ljung-Box sobre sus residuos. \\
		OE4 & LightGBM y N-HiTS producen previsiones para todo el periodo de prueba con la temperatura como covariable, y la memoria recoge los hiperparámetros finales y el espacio de búsqueda explorado. \\
		OE5 & El protocolo queda descrito en el marco metodológico con los orígenes de previsión, las variables exógenas comunes, el horizonte y el presupuesto de búsqueda, y se aplica sin excepción a los cuatro modelos. \\
		OE6 & Existe una tabla con MAE, RMSE, MAPE y MASE de cada modelo sobre 2024 completo y una matriz de $p$-valores del test de Diebold-Mariano para cada par de modelos. \\
		OE7 & Para cada modelo se reportan cobertura empírica, anchura media, pérdida pinball y CRPS de los intervalos al 90\,\% antes y después de aplicar ACI, junto con sus diagramas de fiabilidad. \\
		OE8 & Las métricas de precisión y de calibración se presentan desglosadas por hora del día y por tipo de día, y la discusión identifica los regímenes en los que cambia el modelo preferible. \\
		OE9 & El repositorio público contiene el código, las semillas, las particiones y las dependencias fijadas, y permite reproducir las tablas de resultados de la memoria. \\
		\hline
	\end{tabularx}
\end{table}